Comparando distintos enfoques de selección y reducción de dimensionalidad sobre
de clasificación médica, se lograron contrastar sus principios teóricos con su
desempeño práctico.

La distribución original de 16 clases incluye categorías con muestras de 2 a 5,
lo cual impide una validación significativa. Así, mediante la reformulación
binaria, casi se balancearon los casos. Y de esta manera no se necesitó aplicar
alguna técnica de sobremuestreo o remuestreo.

En validación cruzada, SFS fue superior ($0.778$), pero en prueba Fisher obtuvo
la mejor exactitud ($0.758$) y AUC ($0.764$).


El \textbf{Fisher Score} ofreció un ranking coherente con SFS
en su característica más importante (\texttt{qrs\_duration}).

Por otro lado, \textbf{PCA} no selecciona variables originales, las recombina.
La proyección 2D mostró que las direcciones de máxima varianza no coinciden con
las que discriminan las clases, es decir tiene desempeño predictivo.

La diferencia observada entre la exactitud en validación cruzada y la
exactitud en prueba (especialmente en SFS, $0.778$ vs $0.714$) se debe a la
poca variabilidad de las muestras pequeñas, por ejemplo, $n = 91$.

Finalmente, la eliminación de columnas constantes, imputación, escalado
\textit{z-score} y eliminación de correlaciones para PCA tuvo un impacto
comparable al de la elección del método de selección. Ningún algoritmo compensa
datos mal preparados, y en problemas médicos como este la calidad inicial es
tan determinante como el modelo final.

\QED